\documentclass[11pt]{article}

\usepackage[margin=1in]{geometry}
\usepackage{booktabs}
\usepackage{amsmath}
\usepackage{graphicx}
\usepackage{hyperref}

\title{\vspace{-1.5em}Computational Simulation of Yin's (1969) Face Inversion Effect\\[0.3em]
\large Using a Log-Polar Neural Network and a KDE Memory Model,\\ Averaged Across Four Trained Models}
\author{Siddharth Agrawal}
\date{\today}

\begin{document}
\maketitle

\begin{abstract}
\noindent
We test whether a biologically inspired neural memory model reproduces the classic
result of Yin (1969): that turning a picture upside down hurts memory for faces far more
than it hurts memory for other kinds of objects. We train a log-polar (foveated)
convolutional network from scratch on faces, objects, and houses, read out its internal
binary code, and feed that code into a Kernel Density Estimation (KDE) familiarity model
based on the Barrington NIMBLE framework. We then run Yin's exact study-then-test
recognition design. Averaging over four independently trained models, the network shows a
large inversion cost for faces and only a small one for objects, matching both the pattern
and roughly the size of Yin's human data.
\end{abstract}

\section{Introduction and Background}

Yin (1969) is one of the foundational studies in face recognition. It showed that human
memory for faces is hurt much more by turning the picture upside down than memory for other
familiar, complex objects that are also normally seen in only one orientation
(``mono-oriented''), such as houses, airplanes, and stick figures. This selectivity became
a key piece of evidence that face recognition relies on a special, \emph{holistic}
(whole-at-once) processing mechanism, rather than the generic strategies we use for ordinary
objects.

Yin's task was a forced-choice recognition memory test. Each participant first saw an
\textbf{inspection (study) series} of 40 pictures, shown one at a time for 3 seconds each.
They then saw a \textbf{test series} of 24 pairs of pictures. Every pair contained exactly
one \emph{old} picture (an exact copy of one seen during study) and one \emph{new} picture
(never shown before), and the participant simply pointed to the one they had seen before.
Only 24 of the 40 studied pictures came back in the test, so 16 were never re-tested --- a
realistic memory situation rather than a memorize-everything drill.

Yin ran the study in two orientation setups. \textbf{Experiment I} kept the orientation the
same across study and test: either both upright (Upright--Upright) or both inverted
(Inverted--Inverted). \textbf{Experiment II} flipped the orientation between the two phases:
study upright and test inverted (Up--Down), or study inverted and test upright (Down--Up).
The main finding was that \emph{every} category got harder when inverted, but faces got
disproportionately harder --- pointing to one general ``familiar orientation'' factor that
affects all objects, plus a second, special factor that applies only to faces.

This report reproduces Yin's design computationally. We use a neural network as a stand-in
for the visual system that turns a picture into an internal representation, and a simple
memory model that decides whether a test picture feels ``familiar.'' Two design choices are
worth noting up front: we average results over four separately trained networks instead of
relying on one, and we report the object category alongside faces so the dissociation can be
seen directly.

\section{Converting Yin's Error Scores to Accuracy}

Yin reported performance as the \textbf{mean number of errors} out of 24 test pairs. Our
model outputs a percentage accuracy, so to compare the two we convert Yin's error scores to
accuracy with
\[
\text{Accuracy}(\%) = \frac{24 - M_{\text{errors}}}{24} \times 100 .
\]
Accuracy and error count carry the same information here (one is just $24$ minus the other,
rescaled), so the inversion effect can equally be read as a difference in error counts; we
use accuracy below for a direct side-by-side with the model. Table~\ref{tab:human} gives the
converted human values for the two categories we compare against: faces, and airplanes
(which we treat as the human ``object'' comparison, since our network's object set is closest
in spirit to Yin's simple, unfamiliar object caricatures).

\begin{table}[h]
\centering
\caption{Yin (1969) mean errors (out of 24) with subject-level standard deviations, and the
derived accuracy, for the faces and airplane (``object'') categories across all four orientation
conditions. Error $SD$ values are Yin's originals (Exp.~I Table~1 and Exp.~II Table~2, across
$n=26$ and $n=21$ subjects). The accuracy $SD$ is the error $SD$ rescaled by $100/24$, since
accuracy is a linear transform of error count and so carries the same dispersion.}
\label{tab:human}
\begin{tabular}{llcc}
\toprule
Category & Study $\rightarrow$ Test & Mean Errors ($SD$) & Derived Accuracy ($SD$) \\
\midrule
Faces     & Upright $\rightarrow$ Upright   & 0.89 (1.09) & 96.29\% (4.5) \\
Faces     & Inverted $\rightarrow$ Inverted & 4.35 (1.41) & 81.88\% (5.9) \\
Faces     & Upright $\rightarrow$ Inverted  & 3.81 (1.71) & 84.13\% (7.1) \\
Faces     & Inverted $\rightarrow$ Upright  & 5.14 (1.39) & 78.58\% (5.8) \\
\midrule
Airplanes & Upright $\rightarrow$ Upright   & 3.65 (1.69) & 84.79\% (7.0) \\
Airplanes & Inverted $\rightarrow$ Inverted & 3.85 (2.03) & 83.96\% (8.5) \\
Airplanes & Upright $\rightarrow$ Inverted  & 3.19 (1.94) & 86.71\% (8.1) \\
Airplanes & Inverted $\rightarrow$ Upright  & 4.14 (1.58) & 82.75\% (6.6) \\
\bottomrule
\end{tabular}
\end{table}

Even before any modeling, the human numbers already show the effect: faces drop about
\textbf{15--18 points} when inversion is involved, while airplanes barely move (within about
2 points, and in one condition they even improve).

\section{Computational Simulation: Experimental Design}

\subsection{Model Architecture and Feature Extraction}

The neural model is a convolutional feature extractor in the LPNet style: each fixation is
first foveated (sharp in the center, blurred toward the edges) and then passed through a
\emph{log-polar} transform, which mimics how the retina samples the visual field. Fixation
points themselves are chosen bottom-up from a bank of Gabor filters (a model of V1 simple
cells), so the same, non-face-specific mechanism is used for every category.

The network is trained from scratch (no borrowed pretraining) to classify a combined set of
\textbf{128 face identities, 64 object categories, and houses} (pooled into one generic
``house'' class), all under a single shared classification head. Each training sample is a
$180 \times 180$ patch cropped around one of the Gabor-selected fixation points; before
foveation and the log-polar transform, the patch is randomly rotated by up to $\pm 15^\circ$
(a fresh angle every epoch, modelling small natural variations in head and gaze orientation)
and horizontally flipped with probability $0.5$. The flip is horizontal only: a vertical flip
is never applied during training, since that would leak inverted views into the network and
contaminate the very inversion manipulation the simulation is designed to test. After a bottleneck layer, it
produces a stochastic \textbf{binary hidden representation} $h \in \{0,1\}^{256}$ through a
temperature-scaled sigmoid followed by Bernoulli sampling --- directly analogous to a binary
code over a population of simulated neurons. The classifier head is only used to \emph{train}
the network; the memory model reads the shared 256-dimensional code $h$.

To keep the results from depending on any single trained network, we \textbf{average over
four independently trained models}, spanning two backbone sizes (ResNet-18 and ResNet-34)
and two training fixation counts (16 and 32 fixations per image). Every number reported in
the results tables is a mean across these four models; we never report a single model on its
own.

\subsection{Simulating the Study Phase (Inspection Series)}

Mirroring Yin's 40-picture inspection series, 40 items are loaded for the study phase. For
each item, 10 fixation patches are extracted --- the spots where a viewer's gaze might land
during the brief 3-second exposure --- and passed through the network to get its binary code.
Each item's code is stored in a \textbf{memory bank}, the computational analogue of encoding
a picture into episodic memory.

To add biological realism, \textbf{binomial noise} is applied to each stored code by randomly
flipping bits:
\[
h_{\text{noisy}} = h \oplus m, \qquad m_i \sim \text{Bernoulli}(p_{\text{noise}}).
\]
This models the fact that a memory is not a perfect digital recording but a rough, partly
scrambled impression.

\subsection{Simulating the Test Phase (24 Forced-Choice Pairs)}

Matching Yin's 24-pair test series, 24 \emph{old} items are drawn from the 40 studied ones
and paired with 24 \emph{new} items never seen during study. Each test item is viewed through
32 fixation patches (a richer look, reflecting the slower, self-paced test phase). Noise is
again applied to these retrieval codes, modelling the fact that pulling up a memory is itself
imperfect.

A \textbf{familiarity score} is then computed for each test item using a Kernel Density
Estimation (KDE) rule from the Barrington NIMBLE framework:
\[
p(f \mid c) = \frac{1}{|M_c|} \sum_{m \in M_c}
\exp\!\left(-\frac{\lVert f - m \rVert^2}{2\sigma^2}\right),
\]
where $f$ is the code from a single test fixation, $M_c$ is the set of stored codes for a
studied item $c$, and $\sigma$ is the kernel bandwidth (how forgiving the similarity is). The
familiarity of a whole test item is the best (highest log-likelihood) match against any
stored item. The decision rule follows Yin exactly: \textbf{if the old item scores higher
than the new item, the trial is correct.}

\subsection{Noise Calibration}

The noise level $p_{\text{noise}}$ was calibrated separately for each model and each category
on its own Upright--Upright condition, choosing the smallest noise that brings model accuracy
down to the matching human baseline. For faces the target is Yin's human faces baseline
($\approx 96\%$), and across the four models this landed between $p = 0.00$ and $p = 0.15$
(mean $\approx 0.06$). For objects the target is the human object (airplane) baseline of
$84.79\%$; here the models were already at or below this level with no added noise, so objects
were run at $p = 0.00$. In every case the calibrated noise was then held fixed across all four
orientation conditions for that model.

\subsection{Why the Noise}

Noise enters at two points, and both are intentional. During the \textbf{study phase}, the
model glances at each item through 10 fixations and stores a corrupted version of what it
sees --- the idea that a few seconds of looking leaves a rough impression, not a photograph.
During the \textbf{test phase}, each item is viewed through 32 fixations and noise is applied
again before any comparison --- the separate idea that \emph{retrieving} a memory is itself
imperfect, so even a well-stored trace does not come back fully intact. Using one noise level
at both stages keeps the model simple, and tying its value to the human Upright--Upright
baseline means the noise is grounded in the data rather than chosen arbitrarily.

\section{Results}

\subsection{Inversion Effect During Training}

Before the memory simulation, the inversion effect is already visible in the trained
network's own classification accuracy: we simply measure how well each model classifies
upright pictures versus the same pictures inverted (a 180$^\circ$ rotation applied before the
log-polar transform). Table~\ref{tab:train} shows the mean across the four models.

\begin{table}[h]
\centering
\caption{Classification accuracy on upright vs.\ inverted pictures, averaged across the four
trained models. ``Drop'' is the accuracy lost to inversion.}
\label{tab:train}
\begin{tabular}{lccc}
\toprule
Category & Upright & Inverted & Drop \\
\midrule
Faces   & 93.9\% & 61.2\% & $-32.7$ \\
Houses  & 95.8\% & 74.5\% & $-21.3$ \\
Objects & 77.2\% & 72.0\% & $-5.2$ \\
\bottomrule
\end{tabular}
\end{table}

This is exactly the ordering Yin found. Faces lose about \textbf{33 points} to inversion,
houses about \textbf{21 points}, and objects only about \textbf{5 points}. The face-specific
cost is roughly six times the object cost, with houses falling in between --- a clean
face~$>$~houses~$>$~objects hierarchy that echoes Yin's ``general familiarity factor plus a
special face factor.''

\subsection{Results of the Yin Simulation}

We now report the full Yin 2AFC memory simulation, again as the mean across the four models.
We show \textbf{faces} and \textbf{objects}. (Houses are left out here: because all houses
were trained as a single generic class, the network's code cannot tell one house from
another, so the old/new memory test for houses is at chance and not meaningful. Houses remain
useful for the training-accuracy comparison above, where per-house identity is not required.)

Table~\ref{tab:faces} gives faces and Table~\ref{tab:objects} gives objects. Alongside the
model and human accuracies, we list the \textbf{inversion cost}: how far each condition falls
below its own Upright--Upright baseline. This ``difference between rows'' is the number that
makes the dissociation visible.

Two things should be read carefully alongside these tables. First, \textbf{the Upright--Upright
row is fit, not predicted.} Because $p_{\text{noise}}$ is calibrated on each model's own
Upright--Upright condition to match the human baseline (Section~4.4), the agreement in the UU
row is true by construction and is \emph{not} independent evidence. The genuine model predictions
are the three inversion costs --- the Model~$\Delta$ column --- which are what should be compared
against the human $\Delta$. Second, \textbf{the human ``object'' column is Yin's airplanes.} Our
network's object category is a generic $64$-class object set, whereas Yin's closest mono-oriented
object comparison is airplanes; we therefore compare model objects against human airplanes
throughout. The two are not the same stimuli, so the object comparison is a match of the overall
\emph{pattern} (small, roughly flat inversion cost) rather than a stimulus-for-stimulus fit.
Finally, the model $SD$ is computed across only four trained models and each condition is scored
on $24$ pairs, so the standard deviations are large relative to the object $\Delta$s: the face
inversion cost clears its own $SD$, while the $2$--$5$ point object costs sit within it.

\begin{table}[h]
\centering
\caption{Faces: model accuracy (mean of four models, with the standard deviation across those
four models in parentheses) vs.\ Yin's human accuracy (subject-level $SD$ in parentheses), with
the inversion cost relative to Upright--Upright.}
\label{tab:faces}
\begin{tabular}{llcccc}
\toprule
Study & Test & Model Acc. ($SD$) & Human Acc. ($SD$) & Model $\Delta$ & Human $\Delta$ \\
\midrule
Upright  & Upright  & 95.8\% (0.0) & 96.29\% (4.5) & --- & --- \\
Inverted & Inverted & 80.2\% (5.2) & 81.88\% (5.9) & $-15.6$ & $-14.4$ \\
Upright  & Inverted & 78.1\% (6.3) & 84.13\% (7.1) & $-17.7$ & $-12.2$ \\
Inverted & Upright  & 76.0\% (6.3) & 78.58\% (5.8) & $-19.8$ & $-17.7$ \\
\bottomrule
\end{tabular}
\end{table}

\begin{table}[h]
\centering
\caption{Objects: model accuracy (mean of four models, with the standard deviation across those
four models in parentheses) vs.\ Yin's human \emph{airplane} accuracy (subject-level $SD$ in
parentheses; airplanes are Yin's object category we compare against, see text), with the
inversion cost relative to Upright--Upright.}
\label{tab:objects}
\begin{tabular}{llcccc}
\toprule
Study & Test & Model Acc. ($SD$) & Human Acc. ($SD$) & Model $\Delta$ & Human $\Delta$ \\
\midrule
Upright  & Upright  & 80.2\% (6.3) & 84.79\% (7.0) & --- & --- \\
Inverted & Inverted & 75.0\% (3.4) & 83.96\% (8.5) & $-5.2$ & $-0.8$ \\
Upright  & Inverted & 77.1\% (5.4) & 86.71\% (8.1) & $-3.1$ & $+1.9$ \\
Inverted & Upright  & 78.1\% (4.0) & 82.75\% (6.6) & $-2.1$ & $-2.0$ \\
\bottomrule
\end{tabular}
\end{table}

\subsubsection*{Qualitative fit}

The model reproduces the main qualitative pattern. For \textbf{faces}, Upright--Upright is the
best condition, Upright--Inverted is worse than that and then Inverted--Upright is the worst.\\
For \textbf{objects}, all four conditions sit close together (within about 5 points for the model,
within about 4 for the humans): inversion barely matters. \\
Comparing faces to objects, the accuracy drop when the images are inverted is much worse in faces (16--20pt drop) vs in objects (2--5).

\subsubsection*{Quantitative fit}

The match is not just qualitative. For \textbf{faces}, the inversion costs are close in size: the model
loses $15.6$, $17.7$, and $19.8$ points across the three inverted conditions, against the
human $14.4$, $12.2$, and $17.7$. Upright--Inverted is a weak spot, where the model over-estimates the cost (it makes studying
upright then testing inverted a little harder than humans do), and it slightly swaps the order
of the two middle conditions. But the overall size of the face inversion effect is reproduced,
not just its direction.

For \textbf{objects}, the model's inversion cost stays small ($2$--$5$ points) just as the human effect
does ($0$--$2$ points). The model runs a few points below the human accuracy overall, mainly
because objects are a genuinely harder classification problem for the network, which is why
they already sat at their human baseline with no added noise. Importantly, this works
\emph{against} finding a dissociation --- objects are the harder task here, yet still show the
smaller inversion cost --- so the contrast between faces and objects is if anything
understated.

\subsubsection*{Summary of the dissociation}

Putting the two tables together: inversion ``hurts'' faces by roughly $16$--$20$ points and
``hurts'' objects by only $2$--$5$ points, both in line with the human data. The special cost
that Yin attributed to faces shows up in the model even though nothing about faces was built
into the fixation mechanism or the memory rule --- the same pipeline was used for every
category.

\section{Conclusion}

A from-scratch log-polar network, read out through a simple KDE familiarity memory, reproduces
Yin's (1969) central result: memory for inverted faces is disproportionately impaired compared
with memory for inverted objects. The effect appears twice over --- in the network's raw
classification accuracy (faces lose $\sim$33 points to inversion, objects only $\sim$5) and in
the full study-then-test memory simulation (a $16$--$20$ point face cost against a $2$--$5$
point object cost), with both the pattern and the rough magnitude matching the human data.
Because the model treats faces, objects, and houses with one shared mechanism, the emergence
of a face-selective inversion cost suggests that at least part of Yin's effect can arise from
generic foveated, orientation-sensitive visual processing rather than requiring a hand-built,
face-only module.

\section*{References}

\begin{itemize}
\item Barrington, L., Marks, T.~K., \& Cottrell, G.~W. (2008). NIMBLE: A kernel density model
of saccade-based visual memory. \emph{Journal of Vision, 8}(14), 1--14.
\item Yin, R.~K. (1969). Looking at upside-down faces. \emph{Journal of Experimental
Psychology, 81}(1), 141--145.
\end{itemize}

\end{document}

 